\section{Améliorations et recommandations}
Même si la conception proposée répond aux principales exigences du cahier des charges,
plusieurs améliorations peuvent être envisagées afin d'optimiser le fonctionnement, la fabrication et l'utilisation du robot.
\subsection{Optimisation de la rigidité} Une première amélioration concerne la rigidité globale du système. Les bras du robot sont relativement longs et peuvent être soumis à des efforts de flexion lors des phases d'accélération ou de décélération. 
Une optimisation de leur géométrie permettrait de réduire les déformations tout en limitant l'augmentation de masse.
 Il serait notamment possible d'étudier des formes ajourées ou nervurées.
\subsection{Réduction des jeux mécaniques} 
La précision du robot dépend fortement des jeux présents dans les liaisons mécaniques.
 Une attention particulière devrait donc être portée aux ajustements entre les arbres, les roulements, les paliers et les pièces de fixation. 
 \newline
 L'utilisation de tolérances plus adaptées, de systèmes de précharge ou de liaisons mécaniques plus rigides pourrait permettre de réduire les jeux angulaires. 
Cette amélioration serait particulièrement importante si le robot devait être utilisé pour des démonstrations nécessitant une bonne précision de positionnement.
\subsection{Amélioration esthétique}
Comme le robot est destiné à servir de support de démonstration, son apparence visuelle constitue également un aspect important du projet.
Une amélioration possible serait d'harmoniser la forme des bras, des supports moteurs et des pièces de liaison afin d'obtenir un ensemble plus homogène. L'ajout de caches transparents pour les courroies ou certaines parties de transmission pourrait également améliorer l'aspect général du système, tout en protégeant les éléments mobiles.
\newline
Ces améliorations permettraient d'obtenir un robot plus propre visuellement et plus adapté à une utilisation en démonstration.